\chapter{凸領域}
    \section{凸領域内のベクトルの凸結合は元の凸領域に属す}
        \begin{shadebox}
            $V$を凸領域とし，$\bm{v}_1,\dotsc,\bm{v}_n\in V,\;\alpha_1,\dotsc,\alpha_n\geq 0,\;\sum_{i=1}^n\alpha_i = 1$であるとき，$\sum_{i=1}^n\alpha_i\bm{v}_i \in V$
        \end{shadebox}
        \begin{proof}
            \quad\par
            $n=1$のときは明らか。
            $n=m\in\naturalNumbers$まで成り立つと仮定する。
            $n=m+1$のとき，
            \[ \sum_{i=1}^n\alpha_i\bm{v}_i = \sum_{i=1}^{m+1}\alpha_i\bm{v}_i = \sum_{i=1}^m\alpha_i\bm{v}_i + \alpha_{m+1}\bm{v}_{m+1} = A_m\sum_{i=1}^m\beta_i\bm{v}_i + \alpha_{m+1}\bm{v}_{m+1} \quad \left(A_m \coloneq \sum_{i=1}^m\alpha_i,\;\beta_i \coloneq \frac{\alpha_i}{A_m}\right) \]
            $\sum_{i=1}^m\beta_i = 1$であるから，仮定より$\bm{w}_{m+1} \coloneq \sum_{i=1}^m\beta_i\bm{v}_i\in V$である。
            そして$A_m\geq 0,\;A_m + \alpha_{m+1} = 1$であるから$\sum_{i=1}^n\alpha_i\bm{v}_i = A_m\bm{w}_{m+1} + \alpha_{m+1}\bm{v}_{m+1} \in V$である。
        \end{proof}